\section{Fairness Notions}
\label{sec:notions}

We now describe all the different fairness notions we consider in this paper.

\subsection{Envy-Based Notions}

\begin{definition}[EF]
\label{defn:ef}
Let $\Ical \defeq ([n], [m], (v_i)_{i=1}^n, w)$ be a fair division instance.
In an allocation $A$, an agent $i \in [n]$ \emph{envies} another agent $j \in [n] \setminus \{i\}$ if
$\frac{v_i(A_i)}{w_i} < \frac{v_i(A_j)}{w_j}$.
Agent $i$ is \emph{envy-free} in $A$ (or $A$ is EF-fair to $i$) if
she doesn't envy any other agent in $A$.
\end{definition}

For unequal entitlements, most papers use the term WEF instead of EF.
But we use the term EF in this paper to emphasize that unequal entitlements
is a property of the fair division setting, not the fairness notion.
%
It is easy to see that EF allocations may not exist, so several relaxations have been studied.
Two of the most popular relaxations of EF are
EF1 \cite{budish2011combinatorial,lipton2004approximately},
and EFX \cite{caragiannis2019unreasonable}.

\begin{definition}[EF1]
\label{defn:ef1}
For a fair division instance $\Ical \defeq ([n], [m], (v_i)_{i=1}^n, w)$,
an allocation $A$ is EF1-fair to agent $i$ if for every other agent $j$,
either $i$ does not envy $j$,
or $\displaystyle \frac{v_i(A_i)}{w_i} \ge \frac{v_i(A_j \setminus \{g\})}{w_j}$ for some $g \in A_j$,
or $\displaystyle \frac{v_i(A_i \setminus \{c\})}{w_i} \ge \frac{v_i(A_j)}{w_j}$ for some $c \in A_i$.
\end{definition}

\begin{definition}[EFX]
\label{defn:efx}
For a fair division instance $\Ical \defeq ([n], [m], (v_i)_{i=1}^n, w)$,
an allocation $A$ is EFX-fair to agent $i$ if for each $j \in [n] \setminus \{i\}$,
either $i$ doesn't envy $j$, or both of the following hold:
\begin{enumerate}
\ifx\version\versionIjcai
\item $\displaystyle \frac{v_i(A_i)}{w_i} \ge \frac{\max\left(\left\{
    \begin{array}{l}
    v_i(A_j \setminus S): S \subseteq A_j
    \\\quad \textrm{ and } v_i(S \mid A_i) > 0
    \\\quad \textrm{ and } v_i(S \mid A_j \setminus S) \ge 0
    \end{array}\right\}\right)}{w_j}$.
\else
\item $\displaystyle \frac{v_i(A_i)}{w_i} \ge \frac{\max(\{v_i(A_j \setminus S): S \subseteq A_j
    \textrm{ and } v_i(S \mid A_i) > 0 \textrm{ and } v_i(S \mid A_j \setminus S) \ge 0\})}{w_j}$.
\fi
\ifx\version\versionIjcai
\item $\displaystyle \frac{\min\left(\left\{
    \begin{array}{r}
    v_i(A_i \setminus S): S \subseteq A_i \textrm{ and }
    \\ v_i(S \mid A_i \setminus S) < 0
    \end{array}\right\}\right)}{w_i} \ge \frac{v_i(A_j)}{w_j}$.
\else
\item $\displaystyle \frac{\min\left(\left\{v_i(A_i \setminus S): S \subseteq A_i
    \textrm{ and } v_i(S \mid A_i \setminus S) < 0 \right\}\right)}{w_i} \ge \frac{v_i(A_j)}{w_j}$.
\fi
\end{enumerate}
\end{definition}

\Cref{defn:efx} looks very different from the original definition of EFX
given by \cite{caragiannis2019unreasonable}.
Moreover, an alternative definition of EFX, which we call \EFXZero,
is studied in some works \cite{plaut2020almost,chaudhury2021little,chaudhury2024efx}.
In \cref{sec:notions:efx}, we explain why \cref{defn:efx} makes sense
and why it is better than \EFXZero.
We also show that it is equivalent to the original definition when
valuations are submodular, and it is equivalent to \EFXZero{}
when marginals are (strictly) positive or negative.

\subsection{Proportionality-Based Notions}

\begin{definition}[PROP]
\label{defn:prop}
For a fair division instance $\Ical \defeq ([n], [m], (v_i)_{i=1}^n, w)$,
agent $i$'s \emph{proportional share} is $w_iv_i([m])$.
Allocation $A$ is \emph{proportional} (PROP) if $v_i(A_i) \ge w_iv_i([m])$.
\end{definition}

\begin{definition}[PROP1 \cite{conitzer2017fair}]
\label{defn:prop1}
For a fair division instance $\Ical \defeq ([n], [m], (v_i)_{i=1}^n, w)$,
an allocation $A$ is PROP1-fair to agent $i$ if
either $v_i(A_i) \ge w_iv_i([m])$,
or $v_i(A_i \cup \{g\}) > w_iv_i([m])$ for some $g \in [m] \setminus A_i$,
or $v_i(A_i \setminus \{c\}) > w_iv_i([m])$ for some $c \in A_i$.
\end{definition}

Note that \cref{defn:prop1} uses strict inequalities, whereas most papers do not.
We define it this way to make it a slightly stronger fairness notion.
Also, this nuance doesn't affect most results, except those involving binary valuations.

\begin{definition}[PROPx \cite{aziz2020polynomial,li2022almost}]
\label{defn:propx}
For a fair division instance $\Ical \defeq ([n], [m], (v_i)_{i=1}^n, w)$,
an allocation $A$ is said to be PROPx-fair to agent $i$ iff
either $v_i(A_i) \ge w_iv_i([m])$ or both of these conditions hold:
\begin{tightenum}
\item $v_i(A_i \cup S) > w_iv_i([m])$ for every $S \subseteq [m] \setminus A_i$
    such that $v_i(S \mid A_i) > 0$.
\item $v_i(A_i \setminus S) > w_iv_i([m])$ for every $S \subseteq A_i$
    such that $v_i(S \mid A_i \setminus S) < 0$.
\end{tightenum}
\end{definition}

\Cref{defn:propx} looks different from other well-known definitions of PROPx
\cite{aziz2020polynomial,li2022almost}.
This is because those definitions were for more restricted settings,
and for those settings, we show that our definition is equivalent to theirs
(\cref{sec:notions:propx}).

\begin{definition}[PROPm]
\label{defn:propm}
Let $([n], [m], (v_i)_{i=1}^n, w)$ be a fair division instance.
An allocation $A$ is PROPm-fair to agent $i$ iff
$v_i(A_i) \ge w_iv_i([m])$ or both of these conditions hold:
\begin{tightenum}
\item\label{item:propm:chores}$v_i(A_i \setminus S) > w_iv_i([m])$ for every $S \subseteq A_i$
    such that $v_i(S \mid A_i \setminus S) < 0$.
\item\label{item:propm:goods}For every $j \in [n] \setminus \{i\}$, define
    $\tau_j$ to be 0 if $v_i(S \mid A_i) \le 0$ for all $S \subseteq A_j$,
    and $\min(\{v_i(S \mid A_i): S \subseteq A_j \textrm{ and } v_i(S \mid A_i) > 0\})$ otherwise.
    Define $T \defeq \{\tau_j: j \in [n] \setminus \{i\} \textrm{ and } \tau_j > 0\}$.
    Then either $T = \emptyset$, or $v_i(A_i) + \max(T) > w_iv_i([m])$.
\end{tightenum}
\end{definition}

PROPm was originally defined in \cite{baklanov2021achieving} for additive goods.
That definition has a minor error, which \cref{defn:propm} fixes.
We explain the error in \cref{sec:notions:propm}.
Note that PROPm and PROPx are equivalent for chores.

\subsection{Maximin Share and AnyPrice Share}

\begin{comment}
Share-based notions are those where each agent $i$ has a threshold $\mu_i$,
and an allocation $A$ is fair to agent $i$ if $v_i(A_i) \ge \mu_i$.
We get different fairness notions depending on how we define $\mu_i$.
E.g., PROP is a share-based notion.
(A more precise definition of share-based notions is given in \cite{babaioff2022fair}).
We now describe two more share-based notions, maximin share (MMS) and AnyPrice share (APS).
\end{comment}

\begin{definition}[MMS \cite{budish2011combinatorial}]
\label{defn:mms}
For a finite set $M$, let $\Pi_n(M)$ be the set of all $n$-partitions of $M$.
For a function $f: 2^M \to \mathbb{R}$, define
\[ \MMS_f^n(M) \defeq \max_{P \in \Pi_n(M)} \min_{j=1}^n f(P_j). \]
%
For a fair division instance $\Ical \defeq (N, M, (v_i)_{i \in N}, \eqEnt)$,
agent $i$'s \emph{maximin share} (MMS) is given by $\MMS_{v_i}^{|N|}(M)$.
When the instance $\Ical$ is clear from context, we write $\MMS_i$ instead of $\MMS_{v_i}^{|N|}(M)$.
An allocation $A$ is MMS-fair to agent $i$ if $v_i(A_i) \ge \MMS_{v_i}^{|N|}(M)$.
A partition $P \in \Pi_n(M)$ for which $\min_{j=1}^n v_i(P_j) = \MMS_i$
is called agent $i$'s \emph{MMS partition}.
\end{definition}

For unequal entitlements, there are two well-known extensions of MMS:
weighted MMS (WMMS) \cite{farhadi2019fair}
and pessimistic share (pessShare) \cite{babaioff2023fair}.
See \cref{sec:notions:mms} for their formal definitions.
We focus on WMMS in this paper.

Next, we define AnyPrice Share (APS) \cite{babaioff2023fair}.
For any positive integer $m$, define
$\Delta_m \defeq \{x \in \mathbb{R}^m_{\ge 0}: \sum_{j=1}^m x_j = 1\}$.
For any $x \in \mathbb{R}^m$ and $S \subseteq [m]$, define $x(S) \defeq \sum_{j=1}^m x_j$.

\begin{definition}[APS]
\label{defn:aps}
For a fair division instance $\Ical \defeq ([n], [m], (v_i)_{i=1}^n, w)$,
agent $i$'s AnyPrice Share (APS) is defined as
\[ \APS_i \defeq \min_{p \in \mathbb{R}^m}\;\max_{S \subseteq [m]: p(S) \le w_ip([m])} v_i(S). \]
Here $p$ is called the \emph{price vector}.
A vector $p^* \in \mathbb{R}^m$ is called an \emph{optimal} price vector for agent $i$ if
\[ p^* \in \argmin_{p \in \mathbb{R}^m}\;\max_{S \subseteq [m]: p(S) \le w_ip([m])} v_i(S). \]
\end{definition}

\Cref{defn:aps} is slightly different from the original definition given in \cite{babaioff2023fair}.
However, they assume that all items are goods, and for that special case,
their definition is equivalent to \cref{defn:aps}
(c.f. \cref{thm:aps-optimal-price} in \cref{sec:notions:aps}).

\subsection{Derived Notions}

New fairness notions can be obtained by systematically modifying existing notions.
%
We start with two related concepts, epistemic fairness \cite{aziz2018knowledge,caragiannis2023new}
and minimum fair share \cite{caragiannis2023new}.

\begin{definition}[epistemic fairness]
\label{defn:epistemic}
Let $F$ be a fairness notion.
An allocation $A$ is \emph{epistemic-$F$-fair} to an agent $i$ if
there exists another allocation $B$ that is $F$-fair to agent $i$ and $B_i = A_i$.
$B$ is called agent $i$'s \emph{epistemic-$F$-certificate} for $A$.
\end{definition}

\begin{definition}[minimum fair share]
\label{defn:minfs}
For a fair division instance $\Ical \defeq ([n], [m], (v_i)_{i=1}^n, w)$
and fairness notion $F$, let $\Acal(\Ical, F, i)$ be the set of allocations
that are $F$-fair to agent $i$.
%
Then $A$ is minimum-$F$-share-fair to agent $i$ if there exists
an allocation $B \in \Acal(\Ical, F, i)$ such that $v_i(A_i) \ge v_i(B_i)$.
Then $B$ is called agent $i$'s \emph{minimum-$F$-share-certificate} for $A$.
%
Equivalently, an allocation $A$ is \emph{minimum-$F$-share-fair} to agent $i$ if
$v_i(A_i)$ is at least her minimum-$F$-share, defined as
\[ \minFS(\Ical, F, i) \defeq \min_{X \in \Acal(\Ical, F, i)} v_i(X_i). \]
\end{definition}

We now describe pairwise fairness \cite{caragiannis2019unreasonable}
and groupwise fairness \cite{barman2018groupwise}.

\begin{definition}[restricting]
\label{defn:restricting}
Let $\Ical \defeq (N, M, (v_i)_{i \in N}, w)$ be a fair division instance
and $A$ be an allocation. For a subset $S \subseteq N$ of agents, where $|S| \ge 2$,
let $\restrict(\Ical, A, S)$ be the pair $(\Icalhat, \Ahat)$, where
$\Ahat \defeq (A_j)_{j \in S}$, $\Icalhat \defeq (S, \Mhat, (v_j)_{j \in S}, \what)$,
$\Mhat \defeq \bigcup_{j \in S} A_j$, and $\what_j \defeq w_j / \sum_{k \in S} w_k$.
\end{definition}

\begin{definition}[pairwise fairness]
\label{defn:pairwise}
For a fair division instance $\Ical \defeq (N, M, (v_i)_{i \in N}, w)$,
an allocation $A$ is called \emph{pairwise-$F$-fair} to agent $i$ if
for all $j \in N \setminus \{i\}$, $A^{(j)}$ is $F$-fair to $i$ in the instance $\Ical^{(j)}$,
where $(\Ical^{(j)}, A^{(j)}) \defeq \restrict(\Ical, A, \{i, j\})$.
\end{definition}

\begin{definition}[groupwise fairness]
\label{defn:groupwise}
For a fair division instance $\Ical \defeq (N, M, (v_i)_{i \in N}, w)$,
an allocation $A$ is called \emph{groupwise-$F$-fair} to agent $i$ if
for all $S \subseteq N \setminus \{i\}$, $A^{(S)}$ is $F$-fair to $i$ in the instance $\Ical^{(S)}$,
where $(\Ical^{(S)}, A^{(S)}) \defeq \restrict(\Ical, A, \{i\} \cup S)$.
\end{definition}

In this paper, we consider the following derivative notions:
\begin{tightenum}
\item Epistemic envy-freeness (EEF), epistemic EFX (EEFX), epistemic EF1 (EEF1).
\item Minimum EF share (MEFS), minimum EFX share (MXS), minimum EF1 share (M1S).
\item Pairwise proportionality (PPROP), pairwise MMS (PMMS), pairwise APS (PAPS).
\item Groupwise proportionality (GPROP), groupwise MMS (GMMS), groupwise APS (GAPS).
\end{tightenum}
